% Part: many-valued-logic
% Chapter: syntax-and-semantics
% Section: connectives

\documentclass[../../../include/open-logic-section]{subfiles}

\begin{document}

\olfileid[pt]{mvl}{syn}{con}

\olsection{Linguagens e conectivos}

A lógica proposicional clássica, e muitas outras lógicas, usam uma coleção fixa
de \emph{constantes proposicionais} e \emph{conectivos}. Por
exemplo, usamos os seguintes elementos como primitivos:
\begin{enumerate}
\tagitem{prvFalse}{A constante proposicional para a !!{falsity}~$\lfalse$.}{}
\tagitem{prvTrue}{A constante proposicional para a !!{truth}~$\ltrue$.}{}
\item Os conectivos lógicos:
  \startycommalist
  \iftag{prvNot}{\ycomma $\lnot$ (negação)}{}%
  \iftag{prvAnd}{\ycomma $\land$ (conjunção)}{}%
  \iftag{prvOr}{\ycomma $\lor$ (disjunção)}{}%
  \iftag{prvIf}{\ycomma $\lif$ (!!{conditional})}{}%
  \iftag{prvIff}{\ycomma $\liff$ (!!{biconditional})}{}%
\end{enumerate}
\iftag{defNot,defOr,defAnd,defIf,defIff,defTrue,defFalse,defEx,defAll}{%
Além dos conectivos primitivos acima, também usamos símbolos
definidos como abreviações, tais como
\startycommalist
  \iftag{defNot}{\ycomma $\lnot$ (negação)}{}%
  \iftag{defAnd}{\ycomma $\land$ (conjunção)}{}%
  \iftag{defOr}{\ycomma $\lor$ (disjunção)}{}%
  \iftag{defIf}{\ycomma $\lif$ (!!{conditional})}{}%
  \iftag{defIff}{\ycomma $\liff$ (!!{biconditional})}{}%
  \iftag{defFalse}{\ycomma $\lfalse$ (!!{falsity})}{}%
  \iftag{defTrue}{\ycomma $\ltrue$ (!!{truth})}.}{}

Os mesmos conectivos também são usados em lógicas multivaloradas. Contudo,
muitas vezes é útil incluir, por exemplo, diferentes versões da conjunção
na mesma lógica, o que exige símbolos distintos para manter as
versões separadas. Algumas lógicas multivaloradas também incluem
conectivos sem equivalente na lógica clássica. Portanto, seremos um
pouco mais gerais que o habitual.

\begin{defn}
  Uma \emph{linguagem proposicional} consiste em um conjunto $\Lang L$ de
  \emph{conectivos}. Cada conectivo $\star$ tem uma \emph{aridade}; um
  conectivo de aridade~$n$ é dito \emph{$n$-ário}.
  Conectivos de aridade~$0$ também são chamados \emph{constantes};
  conectivos de aridade~$1$ são chamados \emph{unários}, e conectivos de
  aridade~$2$, \emph{binários}.
\end{defn}

\begin{ex}
  A linguagem padrão da lógica proposicional $\Lang L_0$ consiste nos
  seguintes conectivos (com as aridades correspondentes): 
  $\lfalse$~($0$)
  $\lnot$~($1$),
  $\land$~($2$),
  $\lor$~($2$),
  $\lif$~($2$). A maioria das lógicas que consideraremos usará essa linguagem. Algumas
  lógicas, por tradição e convenção, usam símbolos diferentes para certos
  conectivos. Por exemplo, na lógica do produto, o símbolo da conjunção
  costuma ser $\odot$ em vez de~$\land$. Às vezes é conveniente
  acrescentar um operador novo, por exemplo, o operador de determinação $\triangle$
  ($1$-ário).
\end{ex}

\end{document}
